\documentclass[10pt,a4paper]{amsart}
\usepackage{amsmath,amssymb,amsthm,graphicx}
% Retained minimal preamble for the Stacks Project chapter flat.tex (More on Flatness).
% Verbatim excerpt of the chapter preamble: Xy-pic setup (native preamble lines 14-18) and the native
% theorem declarations with the mathematical macros (native preamble lines 154-217). Copied mechanically;
% nothing added, removed or redefined. The wrapper class and the amsmath/amssymb/amsthm/graphicx base
% packages are supplied by the assembler, so the class/externaldocument part of the native preamble is omitted.
% For commutative diagrams we use Xy-pic
\usepackage[all]{xy}

% We use 2cell for 2-commutative diagrams.
\xyoption{2cell}
% Theorem environments.
%
\theoremstyle{plain}
\newtheorem{theorem}[subsection]{Theorem}
\newtheorem{proposition}[subsection]{Proposition}
\newtheorem{lemma}[subsection]{Lemma}

\theoremstyle{definition}
\newtheorem{definition}[subsection]{Definition}
\newtheorem{example}[subsection]{Example}
\newtheorem{exercise}[subsection]{Exercise}
\newtheorem{situation}[subsection]{Situation}

\theoremstyle{remark}
\newtheorem{remark}[subsection]{Remark}
\newtheorem{remarks}[subsection]{Remarks}

\numberwithin{equation}{subsection}

% Macros
%
\def\lim{\mathop{\mathrm{lim}}\nolimits}
\def\colim{\mathop{\mathrm{colim}}\nolimits}
\def\Spec{\mathop{\mathrm{Spec}}}
\def\Hom{\mathop{\mathrm{Hom}}\nolimits}
\def\Ext{\mathop{\mathrm{Ext}}\nolimits}
\def\SheafHom{\mathop{\mathcal{H}\!\mathit{om}}\nolimits}
\def\SheafExt{\mathop{\mathcal{E}\!\mathit{xt}}\nolimits}
\def\Sch{\mathit{Sch}}
\def\Mor{\mathop{\mathrm{Mor}}\nolimits}
\def\Ob{\mathop{\mathrm{Ob}}\nolimits}
\def\Sh{\mathop{\mathit{Sh}}\nolimits}
\def\NL{\mathop{N\!L}\nolimits}
\def\CH{\mathop{\mathrm{CH}}\nolimits}
\def\proetale{{pro\text{-}\acute{e}tale}}
\def\etale{{\acute{e}tale}}
\def\QCoh{\mathit{QCoh}}
\def\Ker{\mathop{\mathrm{Ker}}}
\def\Im{\mathop{\mathrm{Im}}}
\def\Coker{\mathop{\mathrm{Coker}}}
\def\Coim{\mathop{\mathrm{Coim}}}

% Boxtimes
%
\DeclareMathSymbol{\boxtimes}{\mathbin}{AMSa}{"02}

%
% Macros for moduli stacks/spaces
%
\def\QCohstack{\mathcal{QC}\!\mathit{oh}}
\def\Cohstack{\mathcal{C}\!\mathit{oh}}
\def\Spacesstack{\mathcal{S}\!\mathit{paces}}
\def\Quotfunctor{\mathrm{Quot}}
\def\Hilbfunctor{\mathrm{Hilb}}
\def\Curvesstack{\mathcal{C}\!\mathit{urves}}
\def\Polarizedstack{\mathcal{P}\!\mathit{olarized}}
\def\Complexesstack{\mathcal{C}\!\mathit{omplexes}}
% \Pic is the operator that assigns to X its picard group, usage \Pic(X)
% \Picardstack_{X/B} denotes the Picard stack of X over B
% \Picardfunctor_{X/B} denotes the Picard functor of X over B
\def\Pic{\mathop{\mathrm{Pic}}\nolimits}
\def\Picardstack{\mathcal{P}\!\mathit{ic}}
\def\Picardfunctor{\mathrm{Pic}}
\def\Deformationcategory{\mathcal{D}\!\mathit{ef}}

% --- appended 2026-10-12 (Chief): the project own definitions of the reference/slogan/history
% display environments, copied mechanically from the shared Stacks preamble lines 8-12
% (\usepackage{verbatim} plus the three \newenvironment lines, which the chapter preamble keeps
% verbatim from the project preamble). Nothing added, removed or redefined: these environments
% carry no mathematical content and the project itself typesets them as comments. Without these
% three lines a counted statement whose own text uses them fails with Environment slogan undefined.
\usepackage{verbatim}
\newenvironment{reference}{\comment}{\endcomment}
\newenvironment{slogan}{\comment}{\endcomment}
\newenvironment{history}{\comment}{\endcomment}
\usepackage{cleveref}
\crefname{originalreference}{original reference}{original references}
\Crefname{originalreference}{Original reference}{Original references}
\makeatletter
\newlabel{situation-complex-and-divisor}{{0F8S}{source}{}{sourcepointer-1}{}}
\newlabel{situation-complex-and-divisor@cref}{{[originalreference][0][]0F8S}{[1][0][]source}{}{}{}}
\makeatother
\crefname{originallemma}{original lemma}{original lemmas}
\Crefname{originallemma}{Original lemma}{Original lemmas}
\makeatletter
\newlabel{more-algebra-lemma-eta-vanishing-beta-plus-pre}{{0GSW}{source}{}{sourcepointer-2}{}}
\newlabel{more-algebra-lemma-eta-vanishing-beta-plus-pre@cref}{{[originallemma][0][]0GSW}{[1][0][]source}{}{}{}}
\makeatother
\crefname{originallemma}{original lemma}{original lemmas}
\Crefname{originallemma}{Original lemma}{Original lemmas}
\makeatletter
\newlabel{more-algebra-lemma-eta-vanishing-beta-plus}{{0F81}{source}{}{sourcepointer-3}{}}
\newlabel{more-algebra-lemma-eta-vanishing-beta-plus@cref}{{[originallemma][0][]0F81}{[1][0][]source}{}{}{}}
\makeatother
\crefname{originallemma}{original lemma}{original lemmas}
\Crefname{originallemma}{Original lemma}{Original lemmas}
\makeatletter
\newlabel{more-algebra-lemma-eta-locally-free}{{0F7W}{source}{}{sourcepointer-4}{}}
\newlabel{more-algebra-lemma-eta-locally-free@cref}{{[originallemma][0][]0F7W}{[1][0][]source}{}{}{}}
\makeatother
\crefname{originallemma}{original lemma}{original lemmas}
\Crefname{originallemma}{Original lemma}{Original lemmas}
\makeatletter
\newlabel{cohomology-lemma-eta-cohomology-locally-free}{{0GTA}{source}{}{sourcepointer-5}{}}
\newlabel{cohomology-lemma-eta-cohomology-locally-free@cref}{{[originallemma][0][]0GTA}{[1][0][]source}{}{}{}}
\makeatother
\title{For a good triple there is a finite presentation closed subscheme of the divisor on which the derived twist has locally free cohomology of the expected ranks}
\author{The Stacks Project Authors}
\date{stacks-project:flat.tex:local-snapshot-20261011; GNU FDL 1.2 or later}
\begin{document}
\maketitle
\noindent\textit{Editorial note.} All delivered mathematical text is the author's own, copied from the pinned Stacks Project chapter file flat.tex; no statement or proof step was rewritten, reordered or removed.\par\medskip
\noindent\textit{Editorial note.} The preamble is the chapter's own retained minimal preamble (sources/stacks-flat/retained-preamble.tex), a verbatim excerpt of the native Xy-pic setup and of the native theorem declarations and mathematical macros, followed by the project own definitions of the reference, slogan and history display environments (verbatim package plus three newenvironment lines) that the retained file carries at its end and that the project itself typesets as comments. No macro was invented or redefined, and the wrapper class and the amsmath/amssymb/amsthm/graphicx base packages are supplied by the assembler.\par\medskip
\noindent\textit{Editorial note.} The counted statement and proof contain no cross reference, no citation and no included file, so the delivered body is self contained.\par\medskip
\noindent\textit{Editorial note.} Canonical location of the counted proposition: Stacks tag 0GTZ (sources/stacks-flat/pinned-original-source.tex line 13810).\par\medskip
\noindent\textit{Allowed standard prerequisites (editorial).}\begin{itemize}
\item The ideals J\_i(M\textasciicircum{}bullet, f) and their vanishing properties. More on Algebra, Lemmas labelled more-algebra-lemma-eta-vanishing-beta-plus-pre and more-algebra-lemma-eta-vanishing-beta-plus, cited at proof lines 13837-13842. Define the local ideals cutting out T and identify the map T cap U -> D cap U.
\item Local freeness of the derived twist eta\_f M\textasciicircum{}bullet. More on Algebra, Lemma labelled more-algebra-lemma-eta-locally-free, cited at proof lines 13851-13853. Represents Li\textasciicircum{}*L eta\_I M over T cap U by a complex of locally free C-modules.
\item Local freeness of the cohomology of the twisted complex and the rank comparison. Cohomology, Lemma labelled cohomology-lemma-eta-cohomology-locally-free, cited at proof lines 13855-13857. Gives property (3) of the statement.
\item Good triples and their affine charts. Same chapter, Situation labelled situation-complex-and-divisor, invoked in the statement, with the affine charts of the proof of Lemma labelled lemma-complex-and-divisor-blowup-pre. Provides the cover by good charts on which the construction can be glued.
\end{itemize}\medskip
\ifdefined\hypertarget\hypertarget{sourcepointer-1}{}\fi
\noindent\textit{Source reference.} Original reference 0F8S: sources/stacks-flat/pinned-original-source.tex, source lines 13460--13460. Same-chapter pointer into the author environment situation at source line 13460; invoked at proof line None of the counted proof. It is retained only as a navigation pointer; every load-bearing step is inside the delivered span.\par\medskip
\ifdefined\hypertarget\hypertarget{sourcepointer-2}{}\fi
\noindent\textit{Source reference.} Original lemma 0GSW: sources/stacks-chow/tags.txt, source lines 20946--20946. Cross-chapter standard prerequisite located in the project tag list staged in the local snapshot; invoked at proof line 13840 of the counted proof. It is retained only as a navigation pointer; every load-bearing step is inside the delivered span.\par\medskip
\ifdefined\hypertarget\hypertarget{sourcepointer-3}{}\fi
\noindent\textit{Source reference.} Original lemma 0F81: sources/stacks-chow/tags.txt, source lines 19026--19026. Cross-chapter standard prerequisite located in the project tag list staged in the local snapshot; invoked at proof line 13841 of the counted proof. It is retained only as a navigation pointer; every load-bearing step is inside the delivered span.\par\medskip
\ifdefined\hypertarget\hypertarget{sourcepointer-4}{}\fi
\noindent\textit{Source reference.} Original lemma 0F7W: sources/stacks-chow/tags.txt, source lines 19021--19021. Cross-chapter standard prerequisite located in the project tag list staged in the local snapshot; invoked at proof line 13852 of the counted proof. It is retained only as a navigation pointer; every load-bearing step is inside the delivered span.\par\medskip
\ifdefined\hypertarget\hypertarget{sourcepointer-5}{}\fi
\noindent\textit{Source reference.} Original lemma 0GTA: sources/stacks-chow/tags.txt, source lines 20960--20960. Cross-chapter standard prerequisite located in the project tag list staged in the local snapshot; invoked at proof line 13856 of the counted proof. It is retained only as a navigation pointer; every load-bearing step is inside the delivered span.\par\medskip
\begin{lemma}
\label{lemma-complex-and-divisor-blowup-good-T}
Let $X, D, \mathcal{I}, M$ be as in
Situation \ref{situation-complex-and-divisor}.
If $(X, D, M)$ is a good triple, then there exists a
closed immersion
$$
i : T \longrightarrow D
$$
of finite presentation with the following properties
\begin{enumerate}
\item $T$ scheme theoretically contains $D \cap W$
where $W \subset X$ is the maximal open over which the
cohomology sheaves of $M$ are locally free,
\item the cohomology sheaves of $Li^*L\eta_\mathcal{I}M$
are locally free, and
\item for any point $t \in T$ with image $x = i(t) \in W$ the rank
of $H^i(M)_x$ over $\mathcal{O}_{X, x}$ and the rank
of $H^i(Li^*L\eta_\mathcal{I}M)_t$ over $\mathcal{O}_{T, t}$ agree.
\end{enumerate}
\end{lemma}

\begin{proof}
Let $(U, A, f, M^\bullet)$ be an affine chart for $(X, D, M)$
such that $I_i(M^\bullet, f)$ is a principal ideal for all $i \in \mathbf{Z}$.
Then we define $T \cap U \subset D \cap U$ as the closed subscheme
defined by the ideal
$$
J(M^\bullet, f) = \sum J_i(M^\bullet, f) \subset A/fA
$$
studied in More on Algebra, Lemmas
\ref{more-algebra-lemma-eta-vanishing-beta-plus-pre} and
\ref{more-algebra-lemma-eta-vanishing-beta-plus};
in terms of the second lemma we see that $T \cap U \to D \cap U$
is given by the ring map $A/fA \to C$ studied there.
Since $(X, D, M)$ is a good triple we can cover $X$ by
affine charts of this form and by the first of the two
lemmas, this construction glues. Hence we obtain a closed
subscheme $T \subset D$ which on good affine charts as above
is given by the ideal $J(M^\bullet, f)$. Then properties
(1) and (2) follow from the second lemma. Details omitted.
Small observation to help the reader:
since $\eta_fM^\bullet$ is a complex of locally free modules
by More on Algebra, Lemma \ref{more-algebra-lemma-eta-locally-free}
we see that $Li^*L\eta_\mathcal{I}M|_{T \cap U}$
is represented by the complex $\eta_fM^\bullet \otimes_A C$
of $C$-modules. The statement (3) on ranks follows from
Cohomology, Lemma \ref{cohomology-lemma-eta-cohomology-locally-free}.
\end{proof}

\end{document}
